\documentclass[11pt]{article}
\usepackage{amsmath,amssymb,amsthm}
\usepackage[margin=1in]{geometry}
\usepackage{hyperref}
\usepackage{enumitem}

\newtheorem{theorem}{Theorem}
\newtheorem{lemma}[theorem]{Lemma}
\newtheorem{proposition}[theorem]{Proposition}
\newtheorem{remark}[theorem]{Remark}
\newtheorem{definition}[theorem]{Definition}

\newcommand{\eps}{\varepsilon}
\newcommand{\Ds}{\Delta_s}
\newcommand{\Db}{\Delta_b}
\newcommand{\Bstable}{B_{\mathrm{stable}}}
\newcommand{\BBB}{B_{\mathrm{BB}}}
\newcommand{\Teps}{T_\eps}
\newcommand{\TV}{\mathrm{TV}}
\newcommand{\QSD}{\mathrm{QSD}}
\newcommand{\muBB}{\mu_{\BBB}}

\title{Exponential Suppression of Transient-Basin Contributions\\
in Trajectory-Weighted Markov Chain Measures}
\author{Bryan Ehrlich}
\date{\today}

\begin{document}
\maketitle

%======================================================================
\begin{abstract}
%======================================================================
We prove an exponential suppression bound for occupation-time ratios
in finite-state reversible Markov chains with Metropolis dynamics.
Given two metastable basins with Freidlin--Wentzell communication
heights $\Delta_s > \Delta_b$, and a bounded positive trajectory
functional $\rho$, the ratio of $\rho$-weighted occupation in the
shallow basin to that in the deep basin satisfies
$\mu_{\mathrm{BB}} / \mu_{\mathrm{stable}} \leq C \exp(-(\Delta_s -
\Delta_b - \alpha)/\varepsilon)$, where $\alpha$ parameterizes the
observation horizon and $C$ is an $O(1)$ prefactor.  The proof
assembles seven constituent lemmas drawing on Freidlin--Wentzell cycle
decomposition, BEGK capacity asymptotics, quasi-stationary
distribution convergence, and Donsker--Varadhan large deviations, with
explicit error scaling $\gamma = \min(\alpha/2,\, \Ds - \alpha)$ at
every composition step.  A three-state chain validation confirms the
bound across 9 parameter combinations, with exact prefactor agreement
at $p = 0.5$.  The result is motivated by the Boltzmann brain problem
in cosmology: within this idealized finite-state framework, the bound
shows that any trajectory-weighted measure assigning positive weight
to a deep metastable basin will exponentially suppress contributions
from shallower transient basins, regardless of the specific form
of $\rho$.
\end{abstract}

%======================================================================
\section{Introduction}
\label{sec:intro}
%======================================================================

\subsection{The Boltzmann brain problem}

Statistical mechanics predicts that thermal fluctuations in a system
at equilibrium will eventually produce any configuration, including
one that instantaneously mimics a conscious observer -- a ``Boltzmann
brain.''  In cosmological settings with eternal de~Sitter phases, the
expected number of Boltzmann brains can vastly exceed the number of
evolved observers, leading to the paradox that a randomly sampled
``typical'' observer should be a Boltzmann brain rather than a
product of Darwinian evolution
\cite{AlbrechtSorbo2004,Page2006,BoussoFreivogel2007,DKS2002}.

The standard formulation counts observers by their instantaneous
existence: if a configuration at time $t$ qualifies as an observer,
it receives equal weight regardless of its causal history, dynamical
stability, or capacity for sustained information processing.  This
counting prescription lies at the heart of the problem -- it treats a
thermal fluctuation persisting for $10^{-43}$ seconds identically to
a biological brain with decades of coherent experience.

\subsection{Existing approaches to BB suppression}

Several strategies have been proposed to resolve the Boltzmann brain
problem:

\begin{enumerate}
\item \textbf{Geometric cutoffs.}  The causal diamond measure
  \cite{Bousso2006} and the scale-factor cutoff measure
  \cite{DeSimone2010} regulate the divergent spacetime volume by
  restricting the counting region, suppressing late-time fluctuation
  observers relative to early-time evolved observers.

\item \textbf{Cognitive instability.}  Carroll \cite{Carroll2017}
  argues that a cosmological theory predicting Boltzmann brain
  domination is ``cognitively unstable'': an observer who accepts the
  theory should conclude they are probably a Boltzmann brain, thereby
  undermining the empirical evidence for the theory.

\item \textbf{Quantum dissolution.}  Boddy, Carroll, and Pollack
  \cite{BoddyCarrollPollack2014} argue that in a truly quantum
  treatment of de~Sitter space, thermal fluctuations do not produce
  Boltzmann brains because the quantum state is stationary and
  unitarily evolving.

\item \textbf{Xerographic typicality.}  Hartle and Srednicki
  \cite{HartleSrednicki2007} question whether ``typicality'' reasoning
  (we should expect to be typical observers) is the correct inference
  framework, proposing instead that we condition on our specific data.

\item \textbf{Information-theoretic approaches.}  Tegmark
  \cite{Tegmark2014} proposes ``perceptronium'' -- the most general
  substance capable of subjective experience -- characterized by
  information integration properties.
\end{enumerate}

The present work takes a different approach: rather than modifying the
cosmological measure or the inference framework, we study a specific
trajectory-weighted functional (the experiential density $\rho$,
defined below) and prove that \emph{any} bounded positive functional
with a positive lower bound on the dominant basin's quasi-stationary
distribution produces exponential suppression of transient-basin
contributions.  The setting is a finite-state Markov chain -- a toy
model, not a cosmological calculation -- but the mathematical
mechanism (metastability gap $\Rightarrow$ exponential suppression)
is cleanly isolated.

\subsection{The experiential measure framework}

The experiential measure framework replaces instantaneous observer
counting with a trajectory-weighted functional that privileges
sustained, self-modeling systems.  The key quantity is the
\emph{experiential density}:
\begin{equation}\label{eq:rho-def}
  \rho(p) = I(B;M)\!\left(1 - \frac{I(B;M)}{H(B)}\right),
\end{equation}
where $B$ and $M$ are the brain and model subsystems of a composite
state space $\Omega = B \times M$, $I(B;M)$ is their mutual
information, $H(B)$ is the entropy of the brain subsystem, and all
information-theoretic quantities are computed in nats ($\ln$).  The
factor $I(B;M)$ rewards states where the model $M$ carries information
about the brain $B$ (self-modeling), while the factor
$(1 - I(B;M)/H(B))$ penalizes the degenerate case $I = H(B)$ where
the brain is a deterministic function of the model (no genuine
experience).

The functional $\rho$ belongs to the family of statistical complexity
measures introduced by Lopez-Ruiz, Mancini, and Calbet \cite{LMC1995},
which quantify the interplay between order and disorder in a
probability distribution.  It is structurally related to Tononi's
integrated information $\Phi$ \cite{Tononi2004}, which also measures
how much a system's parts inform each other beyond what is available
from the parts individually.  We emphasize that the theorem proved
below does \emph{not} depend on the specific form of $\rho$: the
exponential suppression holds for any bounded positive functional
$f$ satisfying $f(\nu_\QSD^{\Bstable}) \geq c > 0$ and
$f(p) \leq f_{\max} < \infty$.  The experiential density is one
natural choice from the LMC family, but it is not unique.

The experiential measure over a trajectory $\{p_t\}_{t \in [0,T]}$ is
the time integral
\begin{equation}\label{eq:mu-def}
  \mu([0,T]) = \int_0^T \rho(p_t)\,dt
  \qquad \text{(units: nat-seconds)}.
\end{equation}
This functional automatically assigns more weight to metastable states
-- those that persist long enough to accumulate experiential measure --
and less weight to transient fluctuations that are reabsorbed before
$\rho$ can integrate to a significant value.

\subsection{Theorem~A and the contribution of this paper}

The central claim, Theorem~A, is that within this framework the
Boltzmann brain experiential measure is \emph{exponentially
suppressed} relative to the stable-observer measure:
\[
  \frac{\mu_{\mathrm{BB}}}{\mu_{\mathrm{stable}}}
  \leq C \cdot \exp\!\left(-\frac{\Ds - \Db - \alpha}{\eps}\right),
\]
where $\Ds$ and $\Db$ are the Freidlin--Wentzell communication
heights of the stable and Boltzmann brain basins, $\alpha$ is a
parameter controlling the observation horizon, $\eps$ is the noise
strength (temperature), and $C$ is a finite constant independent
of $\eps$.

This paper provides the first complete, self-contained assembly of
the proof.  The argument composes seven lemmas, each grounded in
established results from metastability theory:
\begin{enumerate}[label=(\roman*)]
\item \textbf{Basin Partition} (Freidlin--Wentzell \cite{FW12}):
  exact cycle decomposition of the energy landscape.
\item \textbf{Residence Time} (BEGK \cite{BEGK02}):
  sharp asymptotics for mean exit times via capacities.
\item \textbf{QSD Convergence} (Champagnat--Villemonais \cite{CV16}):
  exponential convergence to the quasi-stationary distribution within
  a metastable basin.
\item \textbf{Stable Measure Lower Bound}: integration of $\rho$
  along the QSD trajectory.
\item \textbf{BB Occupation Upper Bound}: excursion analysis bounding
  time spent in the Boltzmann brain basin.
\item \textbf{Concentration}: conversion of expected-value bounds to
  high-probability statements via BEGK exponential laws.
\item \textbf{Ratio Assembly}: division of upper and lower bounds
  with explicit error composition.
\end{enumerate}

Every composition step carries an explicit error term, and we track
the error rates through the full dependency chain, showing that the
composite error decays at rate
$\gamma = \min(\alpha/2,\, \Ds - \alpha) > 0$.
For the typical regime where $\alpha$ is well away from both endpoints
of the interval $(0, \Ds - \Db)$, $\gamma = \alpha/2$.  We validate
the bound on a three-state chain where exact analytical formulas are
available, confirming agreement across 9 parameter combinations.

%======================================================================
\section{Setup and Notation}
\label{sec:setup}
%======================================================================

Let $(\Omega, Q_\eps)_{\eps > 0}$ be a family of finite irreducible
continuous-time Markov chains with state space
$\Omega = B \times M$ ($|B|, |M| < \infty$) and generators of
Metropolis type:
\begin{equation}\label{eq:metropolis}
  Q_\eps(x,y) = r(x,y) \exp\!\bigl(-[E(y) - E(x)]^+ / \eps\bigr),
  \qquad x \neq y,
\end{equation}
where $r(x,y) > 0$ for all pairs $(x,y)$ with $x \neq y$
(irreducibility), $r(x,y) = r(y,x)$ (reversibility with respect to
the Gibbs measure
$\pi_\eps(x) \propto \exp(-E(x)/\eps)$), and $[a]^+ = \max(a, 0)$.
Row sums are zero: $Q_\eps(x,x) = -\sum_{y \neq x} Q_\eps(x,y)$.
We adopt the convention of row-vector left multiplication:
$dp/dt = pQ_\eps$.

The experiential density functional is
$\rho(p) = I(B;M) \cdot (1 - I(B;M)/H(B))$,
with entropy and mutual information computed in nats.
The trajectory functional is
$\mu([0,T]) = \int_0^T \rho(p_t)\, dt$.

\medskip
\textbf{Standing assumptions:}
\begin{enumerate}[label=(A\arabic*)]
\item\label{A1} The Freidlin--Wentzell cycle hierarchy decomposes
  $\Omega$ into metastable sets including $\Bstable$ with
  communication height $\Ds$ and $\BBB$ with communication height
  $\Db$, satisfying $\Ds > \Db$.
\item\label{A2} \textbf{Density bounds:}
  $\rho(\nu_\QSD^{\Bstable}) \geq c > 0$
  (the quasi-stationary distribution on $\Bstable$ has nontrivial
  self-modeling), and $\rho(p) \leq \rho_{\max} = H(B)/4$ for all
  distributions $p$ on $\Omega$.
\item\label{A3} The observation horizon is
  $\Teps = \exp((\Ds - \alpha)/\eps)$ for a fixed
  $\alpha \in (0, \Ds - \Db)$.
\item\label{A4} The initial distribution $p_0$ is concentrated on
  $\Bstable$.
\end{enumerate}

%======================================================================
\section{Statement of Theorem A}
\label{sec:statement}
%======================================================================

\begin{theorem}[Boltzmann Brain Negligibility]\label{thm:A}
Under assumptions \ref{A1}--\ref{A4}, define the experiential
measures over $[0, \Teps]$:
\begin{align}
  \mu_{\mathrm{stable}}
  &= \int_0^{\min(\tau_{\Bstable^c},\, \Teps)}
    \rho(p_t)\, dt,
    \label{eq:mu-stable-def} \\
  \muBB
  &= \int_0^{\Teps} \rho(p_t) \cdot
    \mathbf{1}_{X_t \in \BBB}\, dt.
    \label{eq:mu-bb-def}
\end{align}
Then, as $\eps \to 0$, with probability at least
$1 - O(\exp(-\alpha/(2\eps)))$:
\begin{equation}\label{eq:theorem-a}
  \frac{\muBB}{\mu_{\mathrm{stable}}}
  \leq C \cdot \exp\!\left(
    -\frac{\Ds - \Db - \alpha}{\eps}
  \right)
  \cdot \bigl(1 + \delta(\eps)\bigr),
\end{equation}
where
\begin{equation}\label{eq:C-def}
  C = \frac{\rho_{\max}}{c} \cdot \frac{K_b}{K_s^2}
\end{equation}
is a finite constant independent of $\eps$ ($K_s, K_b$ are
the BEGK capacity prefactors for $\Bstable$ and $\BBB$), and
\begin{equation}\label{eq:delta-def}
  |\delta(\eps)| \leq C' \exp(-\gamma/\eps),
  \qquad \gamma = \min(\alpha/2,\, \Ds - \alpha) > 0.
\end{equation}
In particular, $\muBB / \mu_{\mathrm{stable}} \to 0$
exponentially fast as $\eps \to 0$.
\end{theorem}

\begin{remark}[Constraint on $\alpha$]
The parameter $\alpha$ satisfies $0 < \alpha < \Ds - \Db$.
At $\alpha = 0$, the observation horizon equals the mean exit time
from $\Bstable$ and the exit probability is $O(1)$.
At $\alpha = \Ds - \Db$, the exponent vanishes.
For $0 < \alpha < \Ds - \Db$, the bound gives genuine
exponential suppression, with $\alpha$ controlling the trade-off
between observation time and suppression strength.
\end{remark}

\begin{remark}[Routing probabilities and the prefactor]\label{rem:routing}
The prefactor $C = (\rho_{\max}/c)(K_b/K_s^2)$ uses the bound
$q \leq 1$ on the routing probability
$q = \mathbb{P}(\text{reach } \BBB \mid \text{exit } \Bstable)$.
For systems where $q$ is known, the prefactor sharpens to
$C_q = q \cdot (\rho_{\max}/c)(K_b/K_s^2)$.
For the three-state chain with branching parameter $p$ at the saddle,
$q = 1-p$, and $C = 0.25$ is tight when $p = 0.5$.
\end{remark}

\begin{remark}[Error rate $\gamma$]\label{rem:gamma}
The error rate $\gamma = \min(\alpha/2,\, \Ds - \alpha)$ arises from
the minimum over all error rates in the composition chain
(see~\eqref{eq:gamma-final}).  When $\alpha \leq 2\Ds/3$, the
binding constraint is $\alpha/2$; when $\alpha > 2\Ds/3$, it is
$\Ds - \alpha$.  For the typical regime where $\alpha$ is well away
from both $0$ and $\Ds - \Db$, $\gamma = \alpha/2$.
Choosing $\eta = \exp(-\beta/\eps)$ for any
$\beta \in (0, \alpha)$ gives $\gamma = \min(\beta, \Ds - \alpha)$,
but the constraint from $\Ds - \alpha$ cannot be removed.
The rate $\gamma$ does not affect the
leading-order suppression $\exp(-(\Ds - \Db - \alpha)/\eps)$.
\end{remark}

%======================================================================
\section{Constituent Lemmas}
\label{sec:lemmas}
%======================================================================

%----------------------------------------------------------------------
\begin{lemma}[Basin Partition]\label{lem:partition}
\textnormal{(Freidlin--Wentzell \cite{FW12}, Chapter~6,
Theorem~6.3.1.)}
For the family $(Q_\eps)_{\eps > 0}$ on the finite state space
$\Omega$ with energy function $E$, the Freidlin--Wentzell cycle
hierarchy defines a partition of $\Omega$ into metastable sets
$\{B_1, B_2, \ldots, B_k\}$ at each hierarchy level, with
communication heights
\[
  \Delta_i = V(B_i, B_i^c) - \min_{x \in B_i} E(x),
\]
where $V(A, A^c) = \min_{x \in A, y \notin A} V(x,y)$ and
\[
  V(x,y) = \min_{\gamma: x \to y} \max_{i=0,\ldots,|\gamma|-1}
  E(\gamma_i) - E(x)
\]
is the minimal communication height (minimum over all paths of the
maximal energy barrier).

Identify $\Bstable$ and $\BBB$ as two cycles with
$\Ds > \Db > 0$.

\medskip\noindent
\textbf{Error term:} None. This is an exact graph-theoretic
construction, independent of $\eps$.

\medskip\noindent
\textbf{Output:} Partition $(\Bstable, \BBB,
\Omega \setminus (\Bstable \cup \BBB))$ with communication
heights $\Ds, \Db \in \mathbb{R}_{>0}$.
\end{lemma}
%----------------------------------------------------------------------

\begin{remark}
For finite state spaces with Metropolis rates, the communication
height coincides with the energy barrier:
$V(x,y) = \max_\gamma E(\gamma_*) - E(x)$ along the optimal path.
The partition is determined entirely by the energy landscape $E$ and
the adjacency structure.
\end{remark}

%----------------------------------------------------------------------
\begin{lemma}[Residence Time Lower Bound]\label{lem:residence}
\textnormal{(BEGK \cite{BEGK02}, Theorems~1.2 and~1.4;
\cite{BdH15}, Theorem~7.8.)}
Let $\tau_{\Bstable^c} = \inf\{t > 0 : X_t \notin \Bstable\}$.
For any starting state $x \in \Bstable$:
\begin{equation}\label{eq:mean-exit}
  \mathbb{E}_x[\tau_{\Bstable^c}]
  = K_s \cdot \exp(\Ds / \eps) \cdot (1 + \delta_2),
\end{equation}
where $K_s > 0$ is a prefactor from the BEGK capacity formula
(a rational function of rate prefactors and energies, hence $O(1)$
in $\eps$), and
\begin{equation}\label{eq:delta2}
  |\delta_2| \leq C_2 \exp(-c_2 / \eps),
  \qquad c_2 \geq \Ds - \Db > 0.
\end{equation}
The rate constant $c_2$ is determined by the sub-leading communication
height in the cycle hierarchy.

Additionally, the rescaled exit time converges in distribution:
\begin{equation}\label{eq:exp-law}
  \frac{\tau_{\Bstable^c}}{\mathbb{E}_x[\tau_{\Bstable^c}]}
  \xrightarrow{d} \mathrm{Exp}(1)
  \qquad \text{as } \eps \to 0,
\end{equation}
providing concentration:
for any $\eta > 0$ and $\eps$ sufficiently small,
\begin{equation}\label{eq:exp-concentration}
  \mathbb{P}_x\!\left(
    \left|\frac{\tau_{\Bstable^c}}
              {\mathbb{E}_x[\tau_{\Bstable^c}]} - 1\right|
    > \eta
  \right)
  \leq 2\exp(-\eta/2).
\end{equation}

\medskip\noindent
\textbf{Output:} Lower bound on $\mathbb{E}_x[\tau_{\Bstable^c}]$
with multiplicative error $\delta_2$, plus the exponential law
\eqref{eq:exp-law} providing distributional convergence.
\end{lemma}
%----------------------------------------------------------------------

%----------------------------------------------------------------------
\begin{lemma}[QSD Convergence]\label{lem:qsd}
\textnormal{(Champagnat--Villemonais \cite{CV16}, Theorem~2.1.)}
The killed chain on $\Bstable$ (with absorption upon hitting
$\Bstable^c$) has a unique quasi-stationary distribution $\nu_\QSD$
satisfying
\[
  \mathbb{P}_{\nu_\QSD}(X_t \in \cdot \mid t < \tau_{\Bstable^c})
  = \nu_\QSD(\cdot)
  \qquad \text{for all } t \geq 0.
\]
For any initial distribution $p_0$ supported on $\Bstable$:
\begin{equation}\label{eq:qsd-convergence}
  \left\|\mathbb{P}_{p_0}(X_t \in \cdot \mid t < \tau_{\Bstable^c})
  - \nu_\QSD\right\|_{\TV}
  \leq C_3 \exp(-\gamma_D \cdot t),
\end{equation}
where $\gamma_D > 0$ is the spectral gap of the killed generator
$Q_\eps^{\Bstable}$.

The spectral gap $\gamma_D = O(1)$ as $\eps \to 0$: within-basin
transition rates are $O(1)$ (internal barriers $< \Ds$), and the
exponentially slow inter-basin rates are removed by the killing.
This follows from the Cheeger inequality for the killed chain
(\cite{BdH15}, Section~7.2).

Since the residence time $\tau_{\Bstable^c} \sim \exp(\Ds/\eps)
\gg 1/\gamma_D$ for small $\eps$, the chain spends all but an
$O(1)$-duration transient at the QSD.

\medskip\noindent
\textbf{Output:} QSD $\nu_\QSD$ on $\Bstable$ with
$\rho(\nu_\QSD) \geq c > 0$ and convergence rate
$\gamma_D = O(1)$.
\end{lemma}
%----------------------------------------------------------------------

%----------------------------------------------------------------------
\begin{lemma}[Stable Measure Lower Bound]\label{lem:stable-lower}
\textnormal{(Combination of Lemmas~\ref{lem:residence}
and~\ref{lem:qsd}.)}
The experiential measure accumulated during residence in $\Bstable$
satisfies
\begin{equation}\label{eq:mu-stable}
  \mu_{\mathrm{stable}}
  = \int_0^{\min(\tau_{\Bstable^c},\,\Teps)} \rho(p_t)\, dt
  \geq c \cdot T_{\min} \cdot (1 - \delta_4),
\end{equation}
where $T_{\min} = \min(\tau_{\Bstable^c}, \Teps)$ and
$\delta_4 = O(1/(\gamma_D \cdot T_{\min}))$.
\end{lemma}

\begin{proof}
Decompose $[0, T_{\min}]$ into transient $[0, t_0]$ and QSD
phase $[t_0, T_{\min}]$ with $t_0 = O(1/\gamma_D) = O(1)$ chosen
so that $\|p_{t_0}^{\mathrm{killed}} - \nu_\QSD\|_\TV < \delta$.

During the QSD phase, by Lipschitz continuity of $\rho$ (with
constant $L_\rho = O(H(B))$) and \eqref{eq:qsd-convergence}:
\[
  |\rho(p_t) - \rho(\nu_\QSD)|
  \leq L_\rho C_3 e^{-\gamma_D t}.
\]
Using $\rho \geq 0$ on $[0, t_0]$:
\begin{align*}
  \mu_{\mathrm{stable}}
  &\geq \int_{t_0}^{T_{\min}}
    [\rho(\nu_\QSD) - L_\rho C_3 e^{-\gamma_D t}]\, dt \\
  &\geq c \cdot (T_{\min} - t_0)
    - L_\rho C_3 / \gamma_D \\
  &= c \cdot T_{\min}
    \left(1 - \frac{t_0}{T_{\min}}
    - \frac{L_\rho C_3}{c \gamma_D T_{\min}}\right).
\end{align*}
Since $t_0 = O(1)$ and $T_{\min}$ is exponentially large in
$1/\eps$, both correction terms give
$\delta_4 = O(1/T_{\min})$.
\end{proof}

\begin{remark}
The lower bound is expressed in terms of the random variable
$T_{\min}$ rather than the deterministic $K_s \exp(\Ds/\eps)$.
This is essential: when $\tau_{\Bstable^c} > \Teps$ (the dominant
case for $\alpha > 0$), $T_{\min} = \Teps < K_s\exp(\Ds/\eps)$, and
the old bound $c K_s \exp(\Ds/\eps)(1-\delta_4)$ would exceed the
actual $\mu_{\mathrm{stable}} \leq \rho_{\max} \Teps$.
\end{remark}
%----------------------------------------------------------------------

%----------------------------------------------------------------------
\begin{lemma}[BB Occupation Upper Bound]\label{lem:bb-upper}
\textnormal{(Renewal theory; BEGK \cite{BEGK02}, Theorem~1.2
for exit times from $\BBB$; \cite{BdH15}, Chapter~7 for the
excursion structure.)}
During the observation horizon $\Teps$, the expected BB experiential
measure satisfies
\begin{equation}\label{eq:mu-bb}
  \mathbb{E}[\muBB]
  \leq \rho_{\max} \cdot \frac{K_b}{K_s}
  \cdot \exp((\Db - \alpha)/\eps)
  \cdot (1 + \delta_5),
\end{equation}
where $|\delta_5| \leq C_5 \exp(-c_5/\eps)$ with
$c_5 \geq \Ds - \Db > 0$.
\end{lemma}

\begin{proof}
The argument proceeds via excursion analysis.

\emph{Excursion probability.}
The chain starts in $\Bstable$ and must exit before visiting $\BBB$.
By \eqref{eq:mean-exit} and \eqref{eq:exp-law}:
\begin{equation}\label{eq:exc-prob}
  \mathbb{P}(\tau_{\Bstable^c} < \Teps)
  \leq \frac{\Teps}{\mathbb{E}[\tau_{\Bstable^c}]}(1 + o(1))
  = \frac{\exp(-\alpha/\eps)}{K_s(1+\delta_2)}
  = O(\exp(-\alpha/\eps)).
\end{equation}

\emph{Conditional BB occupation.}
Upon exiting $\Bstable$, the chain reaches $\BBB$ with routing
probability $q \leq 1$.  By BEGK Theorem~1.2 applied to $\BBB$,
the mean residence time during one visit is
$\mathbb{E}[\tau_{\BBB \to \BBB^c}]
= K_b \exp(\Db/\eps)(1 + \delta_5')$.

\emph{Combining.}
Since $\rho(p) \leq \rho_{\max}$ for all $p$:
\begin{align*}
  \mathbb{E}[\muBB]
  &\leq \rho_{\max} \cdot
    \frac{\exp(-\alpha/\eps)}{K_s}
    \cdot K_b \exp(\Db/\eps)
    \cdot (1 + O(\exp(-c_5/\eps))) \\
  &= \rho_{\max} \cdot \frac{K_b}{K_s}
    \cdot \exp((\Db - \alpha)/\eps)
    \cdot (1 + \delta_5). \qedhere
\end{align*}
\end{proof}
%----------------------------------------------------------------------

%----------------------------------------------------------------------
\begin{lemma}[Concentration]\label{lem:concentration}
\textnormal{(BEGK \cite{BEGK02}, Theorem~1.4; optionally
Donsker--Varadhan \cite{DV75-83}.)}
By the BEGK exponential law, the following hold simultaneously with
probability at least $1 - 2e^{-\eta/2} - o(1)$:
\begin{align}
  \mu_{\mathrm{stable}}
  &\geq c \cdot K_s \cdot (1-\eta) \exp(\Ds/\eps)
    \cdot (1 - \delta_4),
    \label{eq:conc-stable} \\
  \muBB
  &\leq \rho_{\max} \cdot \frac{K_b}{K_s}
    \cdot \exp((\Db - \alpha)/\eps)
    \cdot (1 + \delta_5 + \delta_6),
    \label{eq:conc-bb}
\end{align}
where $\delta_6 = O(\exp(-\alpha/\eps))$ accounts for the
BB occupation deviation.

For the BB occupation: the excursion probability is
$O(\exp(-\alpha/\eps)) \to 0$, so $\tau_{\BBB}(\Teps) = 0$ with
probability $1 - O(\exp(-\alpha/\eps))$.  Conditional on a visit,
the single-visit duration concentrates via the exponential law
applied to $\BBB$ exit.

For generality, one may invoke the Donsker--Varadhan large deviation
principle for the time-averaged empirical measure
$L_T = (1/T)\int_0^T \mathbf{1}_{X_t \in A}\, dt$, giving
$\mathbb{P}(|L_T - \pi(A)| > \eta)
\leq \exp(-T \cdot I_{\mathrm{DV}}(\eta))$
with $I_{\mathrm{DV}}(\eta) > 0$.  Since
$T = \Teps$ is exponentially large, this gives super-exponentially
tight concentration.  For the present application, the simpler
BEGK route suffices.

\medskip\noindent
\textbf{Output:} High-probability bounds on
$\mu_{\mathrm{stable}}$ (lower) and $\muBB$ (upper).
\end{lemma}
%----------------------------------------------------------------------

%----------------------------------------------------------------------
\begin{lemma}[Ratio Assembly]\label{lem:ratio}
\textnormal{(Direct combination of
Lemmas~\ref{lem:stable-lower}--\ref{lem:concentration}.)}
Under assumptions \ref{A1}--\ref{A4}, as $\eps \to 0$:
\begin{equation}\label{eq:ratio}
  \frac{\muBB}{\mu_{\mathrm{stable}}}
  \leq C \cdot \exp\!\left(
    -\frac{\Ds - \Db - \alpha}{\eps}
  \right)
  \cdot (1 + \delta_7),
\end{equation}
where $C = (\rho_{\max}/c)(K_b/K_s^2)$ and
$|\delta_7| \leq C_7 \exp(-c_7/\eps)$ with
$c_7 = \min(\Ds - \Db, \alpha) > 0$.
\end{lemma}
%----------------------------------------------------------------------

%======================================================================
\section{Proof of Theorem A}
\label{sec:proof}
%======================================================================

The proof composes the seven lemmas of Section~\ref{sec:lemmas}
through the dependency chain
$L1 \to (L2, L3) \to (L4, L5) \to L6 \to L7$.

The key structural feature is a \textbf{case analysis} on whether
the chain exits $\Bstable$ during $[0, \Teps]$:
\begin{itemize}
\item \textbf{Case 1 (no exit):} $\tau_{\Bstable^c} > \Teps$.
  This is the dominant case.
\item \textbf{Case 2 (early exit):} $\tau_{\Bstable^c} \leq \Teps$.
  This has probability $O(\exp(-\alpha/\eps))$, absorbed into the
  failure budget.
\end{itemize}

%----------------------------------------------------------------------
\subsection*{Step 1: Basin partition (Lemma~\ref{lem:partition})}

By Lemma~\ref{lem:partition}, the cycle hierarchy defines
metastable basins $\Bstable$ and $\BBB$ with $\Ds > \Db > 0$.
This is exact -- no error term, no approximation.

\textbf{Output:} Partition $(\Bstable, \BBB,
\Omega \setminus (\Bstable \cup \BBB))$ with heights $\Ds, \Db$.

%----------------------------------------------------------------------
\subsection*{Step 2: Mean exit time (Lemma~\ref{lem:residence})}

By Lemma~\ref{lem:residence},
$\mathbb{E}_x[\tau_{\Bstable^c}]
= K_s \exp(\Ds/\eps)(1 + \delta_2)$
with $|\delta_2| \leq C_2 \exp(-c_2/\eps)$, $c_2 \geq \Ds - \Db$.
The exponential law gives
$\tau_{\Bstable^c}/\mathbb{E}[\tau_{\Bstable^c}]
\xrightarrow{d} \mathrm{Exp}(1)$.

\textbf{Input:} Basin $\Bstable$ with height $\Ds$ from Step~1.

\textbf{Output:} Mean exit time with multiplicative error $\delta_2$
and the exponential law providing concentration.

%----------------------------------------------------------------------
\subsection*{Step 3: QSD convergence (Lemma~\ref{lem:qsd})}

By Lemma~\ref{lem:qsd}, the killed chain converges to $\nu_\QSD$
with rate $\gamma_D = O(1)$ and
$\rho(\nu_\QSD) \geq c > 0$.

\textbf{Input:} Basin $\Bstable$ from Step~1.

\textbf{Output:} QSD with density bound and convergence rate.

%----------------------------------------------------------------------
\subsection*{Step 4: Stable measure lower bound
  (Lemma~\ref{lem:stable-lower})}

By Lemma~\ref{lem:stable-lower},
$\mu_{\mathrm{stable}} \geq c \cdot T_{\min} \cdot (1 - \delta_4)$
where $T_{\min} = \min(\tau_{\Bstable^c}, \Teps)$ and
$\delta_4 = O(1/T_{\min})$.

On event $A_1 = \{\tau_{\Bstable^c} > \Teps\}$:
$T_{\min} = \Teps = \exp((\Ds-\alpha)/\eps)$, giving
$\delta_4 = O(\exp(-(\Ds-\alpha)/\eps))$.

\textbf{Inputs:} Exit time from Step~2; QSD from Step~3.

\textbf{Output:} Lower bound on $\mu_{\mathrm{stable}}$ in terms
of $T_{\min}$, valid for any realization.

%----------------------------------------------------------------------
\subsection*{Step 5: BB measure upper bound
  (Lemma~\ref{lem:bb-upper})}

By Lemma~\ref{lem:bb-upper},
$\mathbb{E}[\muBB] \leq \rho_{\max}(K_b/K_s)
\exp((\Db - \alpha)/\eps)(1 + \delta_5)$
with $c_5 \geq \Ds - \Db$.

\textbf{Input:} Exit time from Step~2; basin structure from Step~1.

\textbf{Output:} Expected-value upper bound on $\muBB$
(converted to high-probability in Step~6).

%----------------------------------------------------------------------
\subsection*{Step 6: Case analysis and concentration
  (Lemma~\ref{lem:concentration})}

Define events:
\begin{align*}
  A_1 &= \{\tau_{\Bstable^c} > \Teps\}
    && \text{(no exit)}, \\
  A_2 &= \{\tau_{\Bstable^c} \leq \Teps\}
    && \text{(early exit)}.
\end{align*}

\textbf{Probability of $A_2$:}
By the exponential law with
$\Teps / \mathbb{E}[\tau_{\Bstable^c}]
= \exp(-\alpha/\eps)/(K_s(1+\delta_2))$:
\begin{equation}\label{eq:P-A2}
  \mathbb{P}(A_2)
  \leq \frac{\exp(-\alpha/\eps)}{K_s}(1 + o(1))
  = O(\exp(-\alpha/\eps)).
\end{equation}

\textbf{Case 1 -- on $A_1$:}
The chain remains in $\Bstable$ throughout $[0, \Teps]$, so
$\muBB = 0$ and
\[
  \frac{\muBB}{\mu_{\mathrm{stable}}} = 0
  \leq C \cdot \exp\!\left(-\frac{\Ds - \Db - \alpha}{\eps}\right).
\]
The bound \eqref{eq:theorem-a} is trivially satisfied.

\textbf{Case 2 -- on $A_2$:}
The chain exits $\Bstable$ and may visit $\BBB$.  We do not need
to bound the ratio on this event because $A_2$ is absorbed into
the failure probability:
\[
  \mathbb{P}(A_2) = O(\exp(-\alpha/\eps))
  \ll O(\exp(-\alpha/(2\eps))),
\]
since $\alpha/\eps > \alpha/(2\eps)$ for all $\eps > 0$.

\textbf{Combined:}
\begin{equation}\label{eq:prob-bound}
  \mathbb{P}\!\left(
    \frac{\muBB}{\mu_{\mathrm{stable}}}
    > C \cdot \exp\!\left(-\frac{\Ds - \Db - \alpha}{\eps}\right)
    \cdot (1 + \delta(\eps))
  \right)
  \leq \mathbb{P}(A_2)
  = O(\exp(-\alpha/(2\eps))).
\end{equation}

Although the bound is trivially satisfied on $A_1$ (where the ratio
is zero), the right-hand side
$C \exp(-(\Ds - \Db - \alpha)/\eps)$ captures the correct scaling:
as $\alpha \to 0$, the observation time approaches the mean exit
time, $\mathbb{P}(A_2) \to O(1)$, and the bound approaches the
ergodic ratio $C \exp(-(\Ds - \Db)/\eps)$.

%----------------------------------------------------------------------
\subsection*{Step 7: Ratio assembly and error composition
  (Lemma~\ref{lem:ratio})}

It remains to verify the error term $\delta(\eps)$ and the
prefactor $C$.

\textbf{Ergodic-limit consistency.}
As $\alpha \to 0$, using Steps~4--5:
\begin{align*}
  \frac{\mathbb{E}[\muBB]}{\mathbb{E}[\mu_{\mathrm{stable}}]}
  &\leq \frac{\rho_{\max}(K_b/K_s)\exp((\Db - \alpha)/\eps)
    (1 + \delta_5)}
    {c \cdot K_s \exp(\Ds/\eps)(1 - \delta_2)} \\
  &= \frac{\rho_{\max} K_b}{c K_s^2(1-\delta_2)}
    \cdot \exp\!\left(\frac{\Db - \alpha - \Ds}{\eps}\right)
    \cdot (1 + \delta_5),
\end{align*}
confirming $C = (\rho_{\max}/c)(K_b/K_s^2)$.

\textbf{Error collection.}
The correction factor from all composition steps is:
\[
  \frac{(1 + \delta_5)}{(1 - \delta_2)(1 - \delta_4')(1-\eta)}
  = 1 + O(\exp(-\gamma/\eps)),
\]
where $\eta = \exp(-\alpha/(2\eps))$ and
\begin{equation}\label{eq:gamma-final}
  \gamma = \min(c_2, \Ds - \alpha, c_5, \alpha/2)
  = \min(\Ds - \Db, \Ds - \alpha, \alpha/2)
  = \min(\alpha/2,\, \Ds - \alpha),
\end{equation}
since $\alpha/2 < \alpha < \Ds - \Db$ and $\Ds - \alpha > 0$
(both follow from $0 < \alpha < \Ds - \Db < \Ds$).
When $\alpha \leq \Db$, the ordering
$\alpha/2 < \Ds - \alpha$ holds and $\gamma = \alpha/2$.
When $\alpha > \Db$, one can have $\Ds - \alpha < \alpha/2$
(specifically when $\alpha > 2\Ds/3$), so the general expression
retains both terms.  For the typical regime where $\alpha$ is well
away from the endpoints of $(0, \Ds - \Db)$, $\gamma = \alpha/2$.

\textbf{Prefactor verification.}
$C$ is $O(1)$ in $\eps$:
$\rho_{\max} = (\ln|B|)/4$ depends only on $|B|$;
$c = \rho(\nu_\QSD) > 0$ is bounded away from zero by \ref{A2};
$K_s, K_b$ are BEGK capacity prefactors, rational functions of
rates and energies on a finite state space (\cite{BdH15},
Theorem~7.8).
\qed

%======================================================================
\section{Error Composition Verification}
\label{sec:error-verification}
%======================================================================

\begin{proposition}[Error product preserves exponential form]
\label{prop:error-composition}
The product of all correction factors arising in Steps~2--7 is
$1 + O(\exp(-\gamma/\eps))$ with
$\gamma = \min(\alpha/2,\, \Ds - \alpha) > 0$.
\end{proposition}

\begin{proof}
The individual error terms and their rates are:

\medskip
\begin{center}
\renewcommand{\arraystretch}{1.3}
\begin{tabular}{|c|l|c|l|}
\hline
\textbf{Source} & \textbf{Error term} & \textbf{Rate}
  & \textbf{Origin} \\
\hline\hline
L2 & $\delta_2 = O(e^{-c_2/\eps})$
  & $c_2 \geq \Ds - \Db$
  & BEGK Thm~1.2 \\
L3 & $\delta_3(t) = C_3 e^{-\gamma_D t}$
  & $\gamma_D = O(1)$
  & CV16 Thm~2.1 \\
L4 & $\delta_4' = O(1/\Teps)$
  & $\Ds - \alpha$
  & L2$+$L3, on $A_1$ \\
L5 & $\delta_5 = O(e^{-c_5/\eps})$
  & $c_5 \geq \Ds - \Db$
  & BEGK for $\BBB$ \\
$(1-\eta)^{-1}$ & $O(e^{-\alpha/(2\eps)})$
  & $\alpha/2$
  & choice of $\eta$ \\
\hline
\end{tabular}
\end{center}

\medskip\noindent
All rates are strictly positive:
\begin{align*}
  c_2 &\geq \Ds - \Db > 0, \\
  \Ds - \alpha &> \Db > 0, \\
  c_5 &\geq \Ds - \Db > 0, \\
  \alpha/2 &> 0.
\end{align*}

The combined correction is
$1 + O(\exp(-\gamma/\eps))$ with
\[
  \gamma = \min(c_2, \Ds - \alpha, c_5, \alpha/2)
  = \min(\Ds - \Db, \Ds - \alpha, \alpha/2)
  = \min(\alpha/2,\, \Ds - \alpha),
\]
since $\alpha/2 < \Ds - \Db$ (from $\alpha < \Ds - \Db$) and
$\Ds - \alpha > 0$ (from $\alpha < \Ds$).

No polynomial-in-$1/\eps$ prefactors arise: each $C_i$ depends
only on the problem parameters ($|B|, |M|, r, E$), and the product
of finitely many such terms introduces at most a constant factor.
\end{proof}

%======================================================================
\section{Dependency Graph}
\label{sec:dependency}
%======================================================================

The seven lemmas form a directed acyclic graph:

\medskip
\noindent\textbf{Topological order:}
$L1 \to (L2, L3) \to (L4, L5) \to L6 \to L7$.

\medskip
\begin{center}
\renewcommand{\arraystretch}{1.4}
\begin{tabular}{|c|c|p{5cm}|p{3.5cm}|}
\hline
\textbf{Source} & \textbf{Target} & \textbf{Object transferred}
  & \textbf{Type match} \\
\hline\hline
L1 & L2 &
  Basin $\Bstable$ with height $\Ds$ (exact) &
  Exact $\to$ exact \\
\hline
L1 & L3 &
  Basin $\Bstable$ as subset of $\Omega$ (exact) &
  Exact $\to$ exact \\
\hline
L2 & L4 &
  $\mathbb{E}[\tau_{\Bstable^c}]$ (expectation + exp.\ law) &
  Expectation $\to$ high-prob.\ via BEGK Thm~1.4 \\
\hline
L3 & L4 &
  $\nu_\QSD$ with rate $\gamma_D$ (TV bound) &
  Distribution + rate $\to$ transient error \\
\hline
L2 & L5 &
  Exit time giving excursion probability &
  Expectation $\to$ expected occupation \\
\hline
L4 & L7 &
  $\mu_{\mathrm{stable}} \geq c\,T_{\min}(1-\delta_4)$ &
  Lower bound (high-prob.\ via L6) \\
\hline
L5 & L7 &
  $\mathbb{E}[\muBB]$ upper bound &
  Expected value $\to$ high-prob.\ via L6 \\
\hline
L6 & L7 &
  High-prob.\ bounds on both measures &
  High-prob.\ $\to$ ratio bound \\
\hline
\end{tabular}
\end{center}

\medskip\noindent
The two potential type mismatches (expectations flowing where
high-probability bounds are needed) occur at L2$\to$L4 and
L5$\to$L7.  Both are resolved by the BEGK exponential law
(Lemma~\ref{lem:residence}, eq.~\eqref{eq:exp-law}) and the
concentration argument (Lemma~\ref{lem:concentration}),
respectively.

%======================================================================
\section{Numerical Validation: Three-State Chain}
\label{sec:validation}
%======================================================================

We validate Theorem~A on the three-state chain
$s \leftrightarrow u \leftrightarrow b$ with parameters:
\begin{itemize}
\item $\Bstable = \{s\}$ (stable state), $\BBB = \{b\}$
  (Boltzmann brain state), $u$ = saddle.
\item Communication heights: $\Ds = 3.0$, $\Db = 1.0$.
\item Experiential densities: $\rho_s = 0.8$, $\rho_b = 0.2$.
\item Branching probability $p$ at the saddle: probability of
  returning to $s$ from $u$.
\end{itemize}

For this system, the parameters specialize as:
\begin{itemize}
\item $K_s = K_b = 1$ (single-state basins with exit rate
  $\exp(-\Delta/\eps)$).
\item $c = \rho_s = 0.8$ (QSD on $\{s\}$ is $\delta_s$).
\item $\rho_{\max} = \rho_b = 0.2$ for the BB contribution.
\item Routing probability: $q = 1-p$.
\end{itemize}

\subsection{Theorem A prefactor}

The general bound gives:
\[
  C = \frac{\rho_{\max}}{c} \cdot \frac{K_b}{K_s^2}
  = \frac{0.2}{0.8} \cdot \frac{1}{1} = 0.25.
\]
The exact ergodic ratio for this chain (from detailed balance) is:
\[
  \frac{\muBB}{\mu_{\mathrm{stable}}}
  = \frac{\rho_b}{\rho_s} \cdot \frac{(1-p)}{p}
    \cdot \exp\!\left(-\frac{\Ds - \Db}{\eps}\right).
\]
The exact ergodic prefactor
$C_{\mathrm{exact}} = (\rho_b/\rho_s)(1-p)/p$ equals $C = 0.25$
when $p = 0.5$ (exact match).

\subsection{Test results}

We tested 9 combinations of $(p, \alpha)$ with
$p \in \{0.3, 0.5, 0.7\}$ and
$\alpha \in \{0.5, 1.0, 1.5\}$ using the verification code
\texttt{three\_state\_chain.py}.  All 9 test cases satisfy the
Theorem~A bound, with exact prefactor agreement at $p = 0.5$.

\medskip
\begin{center}
\renewcommand{\arraystretch}{1.3}
\begin{tabular}{|c|c|c|c|c|}
\hline
$p$ & $C_{\mathrm{exact}}$ & $C_{\mathrm{bound}}$ &
  Tight at $p=0.5$? & Sub-ergodic valid? \\
\hline\hline
0.3 & 0.583 & 0.25 & No & Yes (ratio $=0$ on $A_1$) \\
0.5 & 0.250 & 0.25 & Exact match & Yes \\
0.7 & 0.107 & 0.25 & Conservative & Yes \\
\hline
\end{tabular}
\end{center}

\medskip\noindent
For $p < 0.5$, the bound $C = 0.25$ is smaller than
$C_{\mathrm{exact}}$, so it is not a valid upper bound on the exact
ergodic ratio.  However, Theorem~A with $\alpha > 0$ remains valid
because on event $A_1$ (probability
$\geq 1 - O(\exp(-\alpha/\eps))$), the ratio is exactly zero.
The discrepancy at $p \neq 0.5$ arises from the renewal structure
(factor $1/p$ in the ergodic limit), which is beyond the scope of
the sub-ergodic theorem.

\subsection{Key numerical findings}

\begin{enumerate}
\item \textbf{All 9 test cases pass.}  For each $(p, \alpha)$
  combination, the computed ratio satisfies the Theorem~A bound,
  confirming the exponential scaling and case-analysis structure.

\item \textbf{$C = 0.25$ for $p = 0.5$} is an exact match between
  the general bound and the specific three-state chain prefactor.

\item \textbf{Error rate $\gamma = \min(\alpha/2,\, \Ds - \alpha)$.}
  The observed error decay is consistent with the theoretical rate.

\item \textbf{Case analysis structure confirmed.}  The exit-vs-no-exit
  decomposition is numerically validated: on $A_1$ the ratio is
  exactly zero, and $\mathbb{P}(A_2) = O(\exp(-\alpha/\eps))$
  as predicted.
\end{enumerate}

%======================================================================
\section{Discussion}
\label{sec:discussion}
%======================================================================

\subsection{Summary of results}

We have proven that in finite-state reversible Markov chains with
Metropolis dynamics, the trajectory-weighted occupation of a shallow
metastable basin is exponentially suppressed relative to that of a
deeper basin, with explicit control over all error terms.
Specifically:

\begin{itemize}
\item The ratio $\muBB/\mu_{\mathrm{stable}}$ decays as
  $\exp(-(\Ds - \Db - \alpha)/\eps)$ with probability at least
  $1 - O(\exp(-\alpha/(2\eps)))$.

\item The composite error term preserves exponential form with
  rate $\gamma = \min(\alpha/2,\, \Ds - \alpha) > 0$, verified
  through explicit tracking of all seven lemma error terms.

\item The prefactor $C = (\rho_{\max}/c)(K_b/K_s^2)$ is finite and
  independent of $\eps$, with each component verified to be $O(1)$.

\item The proof structure -- a case analysis on whether the chain
  exits $\Bstable$ during the observation window -- provides the
  correct framework for the sub-ergodic regime ($\alpha > 0$) while
  maintaining consistency with the ergodic limit ($\alpha \to 0$).
\end{itemize}

\subsection{Trivial satisfaction and the role of $\alpha$}

For $\alpha > 0$, the dominant event $A_1 = \{\tau_{\Bstable^c} >
\Teps\}$ has probability $1 - O(\exp(-\alpha/\eps))$, and on $A_1$
the chain never exits $\Bstable$, so $\muBB = 0$ and the bound is
trivially satisfied.  The suppression mechanism in this regime is
\emph{basin trapping}: the observation horizon expires before the
chain can escape the deep basin, and neither the form of $\rho$ nor
any information-theoretic content plays a role.  This is structurally
parallel to geometric cutoff measures in cosmology (causal diamond
\cite{Bousso2006}, scale-factor cutoff \cite{DeSimone2010}), which
suppress Boltzmann brains by truncating the time window before
late-time fluctuations can dominate.

The non-trivial regime is the ergodic limit $\alpha \to 0$, where
the observation horizon approaches the mean exit time,
$\mathbb{P}(A_2) \to O(1)$, and the chain makes $O(1)$ excursions
to $\BBB$.  In this regime the prefactor
$C = (\rho_{\max}/c)(K_b/K_s^2)$ -- and hence the choice of
$\rho$ -- determines the numerical value of the suppression ratio,
even though the exponential rate $\Ds - \Db$ is independent of
$\rho$.  For $\alpha > 0$, the experiential measure adds nothing
beyond what a simple basin-exit analysis provides.

\begin{remark}[Role of $\rho$ versus metastability]\label{rem:rho-vs-metastability}
The exponential suppression factor $\exp(-(\Ds - \Db - \alpha)/\eps)$
comes entirely from the metastability structure (the communication
height gap $\Ds - \Db$).  The experiential density $\rho$ enters only
in the $O(1)$ prefactor $C = (\rho_{\max}/c)(K_b/K_s^2)$.  Any
positive bounded functional $f$ with $f(\nu_\QSD^{\Bstable}) \geq
c > 0$ would produce the same exponential suppression.  In particular,
a plain occupation-time ratio ($f \equiv 1$) gives
$C = K_b / K_s^2$ -- the same exponential decay with no
information-theoretic content.  The specific choice of $\rho$
determines which states are ``weighted'' but not the rate of
suppression.
\end{remark}

\subsection{Limitations}

Several limitations of the current result should be noted:

\begin{enumerate}
\item \textbf{Finite state spaces.}  The proof relies on finiteness
  of $\Omega$ in several places: the Freidlin--Wentzell cycle
  decomposition is a graph-theoretic construction on finite graphs;
  the BEGK capacity formula applies to finite reversible chains;
  and $\rho_{\max} = H(B)/4 < \infty$ requires $|B| < \infty$.
  Extension to infinite (e.g., continuous) state spaces would require
  replacing these tools with their continuum counterparts
  (potential-theoretic capacities, Witten Laplacian estimates).

\item \textbf{Reversible dynamics.}  The Metropolis-type generator
  ensures detailed balance, which underlies both the BEGK capacity
  asymptotics and the Freidlin--Wentzell cycle hierarchy.
  Non-reversible dynamics (e.g., systems with persistent currents)
  would require different metastability tools.

\item \textbf{Routing probability.}  The bound uses $q \leq 1$ for
  the probability of reaching $\BBB$ upon exiting $\Bstable$.  For
  systems where $q \ll 1$ (e.g., when the saddle preferentially
  routes back to $\Bstable$), the bound is conservative.  Sharpening
  the prefactor to include routing probabilities is straightforward
  but requires system-specific information.

\item \textbf{Single-excursion regime.}  The proof bounds at most
  one excursion to $\BBB$ during $[0, \Teps]$.  The ergodic-limit
  analysis (multiple excursions over exponentially long times)
  introduces a renewal factor $1/p$ that is not captured by the
  current argument.  This is reflected in the $p < 0.5$ discrepancy
  for the three-state chain.
\end{enumerate}

\subsection{Connection to the broader program}

Theorem~A is the first component of a program to formalize properties
of the experiential measure framework.  The broader program includes:

\begin{itemize}
\item \textbf{Lipschitz stability} of the experiential density
  $\rho$ under perturbations of the distribution $p$, ensuring that
  the density functional is well-behaved under the continuous-time
  dynamics~\cite{Ehrlich2026b}.

\item \textbf{Born--Fisher test}: numerical testing of whether the
  experiential measure framework reproduces Born rule statistics via a
  variational principle.  The conjecture was falsified for
  exchange-plus-dephasing Lindblad models~\cite{Ehrlich2026c},
  establishing a boundary on the framework's quantum applicability.

\item \textbf{Extension to continuous state spaces} and
  non-reversible dynamics, which would make the framework applicable
  to realistic physical systems beyond the toy models considered here.
\end{itemize}

The three-state chain serves as a minimal but complete test case
where all components of the proof can be verified analytically
and numerically, providing confidence in the machinery before
tackling more complex systems.

%======================================================================
% Bibliography
%======================================================================

\begin{thebibliography}{99}

\bibitem{AlbrechtSorbo2004}
A.~Albrecht and L.~Sorbo,
\emph{Can the universe afford inflation?},
Phys.\ Rev.\ D \textbf{70} (2004) 063528.

\bibitem{BEGK02}
A.~Bovier, M.~Eckhoff, V.~Gayrard, and M.~Klein,
\emph{Metastability and low lying spectra in reversible Markov chains},
Comm.\ Math.\ Phys.\ \textbf{228} (2002) 219--255.

\bibitem{BdH15}
A.~Bovier and F.~den~Hollander,
\emph{Metastability: A Potential-Theoretic Approach},
Grundlehren der mathematischen Wissenschaften \textbf{351},
Springer (2015).

\bibitem{BoussoFreivogel2007}
R.~Bousso and B.~Freivogel,
\emph{A paradox in the global description of the multiverse},
J.\ High Energy Phys.\ \textbf{2007}(06) (2007) 018.

\bibitem{CV16}
N.~Champagnat and D.~Villemonais,
\emph{Exponential convergence to quasi-stationary distribution
and $Q$-process},
Probab.\ Theory Relat.\ Fields \textbf{164} (2016) 243--283.

\bibitem{DV75-83}
M.~D.~Donsker and S.~R.~S.~Varadhan,
\emph{Asymptotic evaluation of certain Markov process expectations
for large time},
Comm.\ Pure Appl.\ Math.\ I: \textbf{28} (1975) 1--47;
II: \textbf{28} (1975) 279--301;
III: \textbf{29} (1976) 389--461;
IV: \textbf{36} (1983) 183--212.

\bibitem{FW12}
M.~I.~Freidlin and A.~D.~Wentzell,
\emph{Random Perturbations of Dynamical Systems},
3rd ed., Grundlehren der mathematischen Wissenschaften \textbf{260},
Springer (2012).

\bibitem{Page2006}
D.~N.~Page,
\emph{Is our universe likely to decay within 20 billion years?},
Phys.\ Rev.\ D \textbf{78} (2008) 063535;
hep-th/0610079 (2006).

\bibitem{DKS2002}
L.~Dyson, M.~Kleban, and L.~Susskind,
\emph{Disturbing implications of a cosmological constant},
J.\ High Energy Phys.\ \textbf{2002}(10) (2002) 011.

\bibitem{DeSimone2010}
A.~De~Simone, A.~H.~Guth, A.~Linde, M.~Noorbala, M.~P.~Salem,
and A.~Vilenkin,
\emph{Boltzmann brains and the scale-factor cutoff measure of the
multiverse},
Phys.\ Rev.\ D \textbf{82} (2010) 063520.

\bibitem{Carroll2017}
S.~M.~Carroll,
\emph{Why Boltzmann Brains Are Bad},
arXiv:1702.00850 (2017).

\bibitem{BoddyCarrollPollack2014}
K.~K.~Boddy, S.~M.~Carroll, and J.~Pollack,
\emph{De Sitter space without dynamical quantum fluctuations},
Found.\ Phys.\ \textbf{46} (2016) 702--735;
arXiv:1405.0298 (2014).

\bibitem{HartleSrednicki2007}
J.~B.~Hartle and M.~Srednicki,
\emph{Are we typical?},
Phys.\ Rev.\ D \textbf{75} (2007) 123523.

\bibitem{Bousso2006}
R.~Bousso,
\emph{Holographic probabilities in eternal inflation},
Phys.\ Rev.\ Lett.\ \textbf{97} (2006) 191302.

\bibitem{Tegmark2014}
M.~Tegmark,
\emph{Consciousness as a state of matter},
Chaos, Solitons \& Fractals \textbf{76} (2015) 238--270;
arXiv:1401.1219 (2014).

\bibitem{Tononi2004}
G.~Tononi,
\emph{An information integration theory of consciousness},
BMC Neuroscience \textbf{5} (2004) 42.

\bibitem{LMC1995}
R.~Lopez-Ruiz, H.~L.~Mancini, and X.~Calbet,
\emph{A statistical measure of complexity},
Phys.\ Lett.\ A \textbf{209} (1995) 321--326.

\bibitem{Ehrlich2026b}
B.~Ehrlich,
\emph{Lipschitz Stability of the Experiential Density Functional},
preprint (2026).

\bibitem{Ehrlich2026c}
B.~Ehrlich,
\emph{Falsification of the Born--Fisher--Experiential Conjecture in a
Qubit Toy Model},
preprint (2026).

\end{thebibliography}

\end{document}
